\chapter{Random multiplicative functions}\label{ch:oc-multiplicative}

This chapter introduces the Steinhaus and Rademacher models, and the conjectures of Chowla and Sarnak that ask the M\"obius function to behave like them.

\section{Steinhaus and Rademacher models}

In the Steinhaus model, $f(p)$ is uniform on the unit circle and $f$ is completely multiplicative. The values are orthogonal, $\mathbb E\,f(n)\overline{f(m)}=\mathbf 1_{n=m}$, so $\mathbb E|S(x)|^2=\lfloor x\rfloor$ for $S(x)=\sum_{n\le x}f(n)$. Wintner's theorem says that $\sum f(n)/n$ converges almost surely, and the fourth moment has the size $\mathbb E|S(x)|^4\sim\frac{12}{\pi^2}x^2\log x$. Given the values on primes up to $\sqrt x$, the contribution of the larger primes is approximately Gaussian, with a random variance. The Rademacher model is the same construction with random signs $f(p)=\pm1$, supported on squarefree integers.

\section{The M\"obius function}

The M\"obius function $\mu$ is multiplicative and takes values in $\{-1,0,1\}$. Random models are a way to guess its partial sums and its correlations; they are not a substitute for a theorem about $\mu$.

\section{Chowla and Sarnak}

Chowla's conjecture asserts that the autocorrelations of $\mu$ along linear forms tend to $0$. Sarnak's conjecture asserts that $\sum_{n\le x}\mu(n)a_n=o(x)$ for every deterministic sequence $(a_n)$ of zero topological entropy. Both ask $\mu$ to look like a random sign.
